\subsection{I-44. Torus-chain connectors and the next balanced degree windows}
\label{I-44}

\paragraph{Statement.}
Fix arbitrary distinct nodes $a_1,\ldots,a_n$ over a field
$E$, and use the balanced rational two-pole map of I-43.
The following assertions hold in an unknown witnessing
ROABP order.

First, let $f$ be irreducible, homogeneous, multilinear,
common-translation invariant, and of width at most three,
and remove its unused variables. Construct the actual minimal
affine cores from its coefficient support spaces. If every
run of rank-three cores has length at most five, the common
one-pole center $\alpha_i=t+z/(a_i-y)$ preserves all strict
internal interval ranks. This assertion holds in every
characteristic. Consequently balanced two-pole evaluation
hits $f$ whenever those runs have length at most four, or
at most five if $\operatorname{char}E\ne2$.

Second, without an irreducibility or run promise, every
nonzero homogeneous multilinear common-translation-invariant
width-at-most-three polynomial is hit when
\[
 0\le n-2\deg f\le4,
 \quad\hbox{or when}\quad
 0\le n-2\deg f\le5\ \hbox{and }\operatorname{char}E\ne2.
 \tag{I.172}
\]
Together with P-19, this proves the balanced two-pole
conclusion for all nonzero multilinear width-at-most-three
polynomials on at most twelve variables in every
characteristic, and on at most thirteen variables outside
characteristic two.

All the scalar conclusions have the same explicit finite
grid as I-40:
\[
 (n+1)^2\binom{D+3}{3},\qquad
 |E|\ge\max\{2n+1,D+1\},
 \tag{I.173}
\]
for total degree at most $D$, with evaluation over a stated
working extension if necessary. These are conclusions for
the balanced rational map with common slope one. No ordinary
centered-$\mathrm{SV}_2$ or arbitrary-slope promotion is made.

\paragraph{1. Torus spectra that need no anchor classification.}
Use I-40's actual minimal cores. Normalize the constant
matrices in any full-rank run to the identity. If the initial
three integer torus weights are $\lambda_1,\lambda_2,
\lambda_3$, the successive weight lists are obtained by
adding one to every entry. Each linear matrix $N_i$ has an
entry only from a column weight to a row weight two larger.

Unless the three weights are distinct and have the form
\[
 (q+2,q,q-2),
 \tag{I.174}
\]
all the allowed entries have a common one-dimensional image
or a common one-dimensional domain, or there is at most one
allowed entry. Indeed, with three distinct weights two
composable steps require exactly (I.174); without such a
chain there is at most one step. If a weight is repeated,
the only nonzero case has two levels of multiplicities one
and two. All products $N_iN_j$ then vanish, and their linear
span consists of matrices of rank at most one.

It follows that the run product is $I+\sum_i x_iN_i$ and
its second compound is affine in all run variables. After
multiplication by the two anchor compounds it has the form
\[
 P_0(u,v)+\sum_i x_iP_i(u,v).
 \tag{I.175}
\]
The original connector has rank two, so this polynomial
matrix is nonzero. The unique-internal-pole argument of
I-43, part~4, applies without a restriction on the number of
summands: at the node of $x_i$, the two anchor values are
independent linear forms in $t,z$. A nonzero $P_i$ therefore
has nonzero residue; if only $P_0$ remains, pair dominance
applies. This proves preservation for every spectrum other
than (I.174), in every characteristic. The differences of
weights here are integer characters of the torus, not
integers reduced modulo the characteristic.

\paragraph{2. All minimal singular anchors into a three-weight chain.}
Assume (I.174), in the displayed order. A preceding singular
anchor $A(u)$ has three columns and two or three rows. Both
its row and column coefficient tensors are independent over
$E$, by actual minimality, and its generic matrix rank is
two. Torus covariance permits preceding row weights only
from
\[
 q+3,\ q+1,\ q-1,\ q-3.
 \tag{I.176}
\]
Its columns have respectively the following possible
coefficient vectors in those row-weight spaces:
\[
 (u a,b,0,0)^{\mathsf T},\quad
 (0,u c,d,0)^{\mathsf T},\quad
 (0,0,u e,f)^{\mathsf T}.
 \tag{I.177}
\]

No preceding weight can have multiplicity two. At either
extreme of (I.176) all row coefficients use a single vector,
contrary to row independence. At weight $q+1$, multiplicity
two requires $b,c$ independent. A nonzero third column must
then use one of the two lower weights, and the resulting
three-column minor is nonzero, contrary to rank two.
The argument at weight $q-1$ is the same in reverse. Three
copies of a weight cannot be spanned by its at most two
coefficient vectors either.

Thus the preceding weights are distinct. The complete list
of possibilities, with rows ordered by decreasing weight,
is as follows. Every displayed scalar is nonzero, except
that in the middle two-row case $c,d$ may vanish individually
but not both.
\[
\begin{array}{c|c}
 \hbox{preceding weights}&A(u)\\ \hline
 (q+3,q-1)&
 \begin{pmatrix}a u&0&0\\0&d&e u\end{pmatrix}\\[3pt]
 (q+1,q-1)&
 \begin{pmatrix}b&c u&0\\0&d&e u\end{pmatrix}\\[3pt]
 (q+1,q-3)&
 \begin{pmatrix}b&c u&0\\0&0&f\end{pmatrix}\\[3pt]
 (q+3,q+1,q-1)&
 \begin{pmatrix}a u&0&0\\b&0&0\\0&d&e u\end{pmatrix}\\[3pt]
 (q+1,q-1,q-3)&
 \begin{pmatrix}b&c u&0\\0&0&e u\\0&0&f\end{pmatrix}
\end{array}
 \tag{I.178}
\]
To verify completeness, the other two-row subsets miss an
entire column. Among the four three-row subsets, the two
omitted subsets give a nonzero triangular determinant.
For the two retained three-row subsets, rank two forces
the middle diagonal entry to vanish; row and column
independence force all the entries shown as nonzero.
This proves the list directly from (I.177).

Over $E(u)$ the row spaces in (I.178) have one of three
forms:
\[
\begin{aligned}
 G_1(u)&=\begin{pmatrix}b&c u&0\\0&d&e u\end{pmatrix},
       &&bce\ne0,\quad d\in E,\\
 G_2(u)&=\begin{pmatrix}b&c u&0\\0&0&1\end{pmatrix},
       &&bc\ne0,\\
 B(u)&=\begin{pmatrix}1&0&0\\0&d&e u\end{pmatrix},
       &&de\ne0.
\end{aligned}
 \tag{I.179}
\]
The passage to these row spaces uses only full-rank
univariate factors and nonzero univariate denominators.
All remain valid after the one-pole substitution, since
$\alpha_i$ is transcendental over $E$.

\paragraph{3. Either of the first two anchor forms suffices.}
The normalized run product in the chain case is
\[
 M=\begin{pmatrix}1&A&B_0\\0&1&C\\0&0&1\end{pmatrix},
 \quad
 A=\sum_i p_ix_i,\quad C=\sum_i q_ix_i,\quad
 B_0=\sum_{i<j}p_iq_jx_ix_j.
 \tag{I.180}
\]
Here each individual core is
$I+x_i(p_iE_{12}+q_iE_{23})$; no quadratic term occurs
within a single core.

Let the opposite, unreduced minimal anchor have rows
$D_1(v),D_2(v),D_3(v)$. They are independent affine
coefficient tensors, and the matrix $D(v)$ has generic
rank two. If the left row space is $G_1$, the coefficient
of $u^2$ in the exterior row of $G_1MD$ is
$ce\,D_2\wedge D_3$. If this is nonzero, its unique double
pole preserves the connector. If $D_2\wedge D_3=0$,
I-43's primitive affine factorization gives
$D_2=p_2(v)q$, $D_3=p_3(v)q$, with $q$ constant and
$D_1\wedge q\ne0$. The full exterior row becomes
\[
 b\,[d p_2+(dC+e u)p_3]\,(D_1\wedge q).
 \tag{I.181}
\]
Its simple pole at the left-anchor node is nonzero.

For $G_2$ the exterior row is
\[
 b\,D_1\wedge D_3+(bA+c u)D_2\wedge D_3.
 \tag{I.182}
\]
If $D_2\wedge D_3\ne0$, use its unique simple pole in
$u$. Otherwise the expression is the nonzero univariate
vector $bD_1(v)\wedge D_3(v)$, which also survives.
These arguments do not bound the length of the middle run.
Order reversal supplies the same conclusion if the opposite
anchor has either of these two forms in the reversed word.

\paragraph{4. Two remaining anchor forms and a quadratic pole test.}
It remains to consider form $B$ at both ends. At the end
of a run of length $\ell$, the left weights are
$(q+\ell+2,q+\ell,q+\ell-2)$. Homogeneity pairs them
antidiagonally, with nonzero constants, with the three
distinct right weights. Apply (I.179) to the reverse anchor
and remove its full-rank univariate outside factor. After
constant column rescaling and ordering, the actual right
factor in the two-dimensional contraction is
\[
 \begin{pmatrix}\alpha v&0\\\beta&0\\0&1\end{pmatrix},
 \qquad \alpha\beta\ne0.
 \tag{I.183}
\]
This is the actual coefficient pairing; it does not assume
that a three-dimensional module is self-dual.

The two-by-two contraction is therefore
\[
 \begin{pmatrix}
  \alpha v+\beta A&B_0\\
  d\beta&dC+e u
 \end{pmatrix},
 \qquad
 F=(\alpha v+\beta A)(dC+e u)-\beta dB_0,
 \tag{I.184}
\]
where $\alpha,\beta,d,e$ are nonzero. In particular $F$ is
nonzero: its $uv$ coefficient is $\alpha e$.

We record the elementary quadratic restriction imposed by
one-pole annihilation. For any homogeneous quadratic in
distinct coordinate variables, a nonzero square coefficient
has a unique double pole at that coordinate's node.
Consequently an annihilating quadratic must be multilinear.
Write it as $\sum_{i<j}c_{ij}X_iX_j$, with symmetric
coefficients off the diagonal. The coefficient of $tz$ in
its one-pole pullback, and then the simple pole of its $z^2$
coefficient at $y=a_i$, give respectively
\[
 \sum_{j\ne i}c_{ij}=0,
 \qquad
 \sum_{j\ne i}\frac{c_{ij}}{a_j-a_i}=0.
 \tag{I.185}
\]
Two distinct possible neighbors give an invertible two-row
matrix in these equations. Thus every nonzero row of an
annihilating quadratic has at least three nonzero neighbors.
This derivation uses no factor of two and holds in every
characteristic.

Apply this to (I.184). The coefficient of $x_i^2$ is
$\beta d p_iq_i$. If any is nonzero, the connector already
survives. Otherwise define the disjoint sets
$P=\{i:p_i\ne0\}$ and $Q=\{i:q_i\ne0\}$.
In the resulting multilinear quadratic, a row indexed by
$i\in P$ has precisely its nonzero edge to $u$ and its
edges to $Q$ positions strictly before $i$. A row indexed
by $i\in Q$ has its edge to $v$ and its edges to $P$
positions strictly after $i$. This follows from
\[
 AC-B_0=\sum_{i>j}p_iq_jx_ix_j
 \quad\hbox{when all }p_iq_i=0.
 \tag{I.186}
\]

Suppose $\ell\le5$. Then at least one of $|P|,|Q|$ is
at most two. If either set is empty, $B_0=0$ and $F$ is
a product of two nonzero linear forms. Such a product
survives because distinct simple poles make the one-pole
map nonzero on every nonzero linear form.

If both sets are nonempty and $|Q|\le2$, (I.185) applied
to any $P$ row forces $|Q|=2$ and forces both $Q$ positions
to precede that row. Hence every $Q$ precedes every $P$.
This again makes $B_0=0$, so the same nonzero product
survives. If $|P|\le2$, apply the symmetric argument to
the $Q$ rows; again every $Q$ must precede every $P$.
Thus (I.184) cannot annihilate the one-pole map for
$\ell\le5$.

This proves that every constrained rank-two connector with
at most five intervening full-rank cores preserves rank.
All coefficient identities, univariate factors and residues
are in the same field $E(t,z,y)$.

\paragraph{5. Completion of the rank and scalar assertions.}
The irreducible/minimal reduction ensures that strict
internal interval ranks are two or three and that the first
and last internal factors are singular, exactly as in I-41.
The two cases of part~1 and parts~2--4 cover every torus
spectrum and every connector of length at most five.
Univariate pivots remain nonzero. Their evaluated rank
factorizations and the invertible two-by-two adjacent
contractions therefore preserve all strict interval ranks,
as proved in I-43, part~5. Pair dominance supplies the
nonzero active endpoint projections. All these assertions
hold simultaneously with the same one-pole parameters.
Small words with no such internal interval are directly
covered by pair or triple dominance.

I-42, part~3, now proves scalar nonvanishing when all runs
have length at most four, or at most five outside
characteristic two. This proves the first assertion.

For (I.172), factor a hypothetical annihilator and recompute
the chosen factor's actual minimal cores, as in I-41.
Each factor is homogeneous and translation-invariant, uses
a disjoint set of variables, and has nonnegative weight.
Its weight is bounded by the total weight. I-40 bounds its
full-rank runs by that weight, so the first assertion applies
to the required factors. The target of the evaluation map
is a field, so no product of these nonzero images can vanish.
This proves (I.172), without assuming that a run promise
for one nonminimal representation passes to its factors.

\paragraph{6. Twelve and thirteen variables.}
I-43 already handles at most eleven variables. On twelve
variables, P-19 and multilinear reversal leave only
homogeneous degrees four through eight. A quartic
annihilator is translation-invariant: a highest nonzero
positive translation coefficient would inherit width at
most three and have degree at most three, contradicting
P-19 and the known minimum degree. Its weight is four,
contradicting (I.172). Reversal excludes degree eight.
For degree five the same leading-translation argument
uses the already excluded lower degrees and produces
translation invariance of weight two; I-41 excludes it,
as well as degree seven by reversal. Degree six similarly
becomes translation-invariant of weight zero and is
excluded by I-40.

On thirteen variables, outside characteristic two, the
same induction starts with a quartic of weight five,
excluded by (I.172). A quintic becomes translation-invariant
of weight three and is excluded by I-43; a sextic has
weight one and is excluded by I-40. Reversal handles their
complementary degrees. In both cases a lowest nonzero
homogeneous component inherits the witnessing width,
so these homogeneous conclusions cover every multilinear
input. The interpolation proof of I-40 gives (I.173)
without a new degree estimate.\hfill$\square$

\paragraph{Remaining scope.}
The rank-preservation conclusion for five-core runs is
valid even in characteristic two. The scalar conclusion
there remains restricted to four-core runs because the
seven-node abstract Hessian of I-42 shows that its three
kernel equations alone are insufficient. In particular,
the weight-five characteristic-two case is reduced to its
possible saturated length-five runs, with their interval
ranks now preserved; no abstract Hessian is asserted to
integrate to a counterexample.

For unrestricted run length, the two remaining anchor forms
lead to the explicit quadratic (I.184). When both internal
types have at least three positions, the short-run argument
above no longer suffices. This is a constrained connector
problem, not a claim that general shared-parameter matrices
compose. The higher-degree and growing-width sampling gates
remain open. No historical priority claim or unperformed
numerical verification is made.
